\chapter{Deployment Roadmap and Conclusion}
\label{ch:ch08_conclusion}

% --- Merged from ch22_what_this_thesis_proves.tex ---

\section{What This Thesis Proves and What It Leaves Open}
\label{ch:what-proves}

A dissertation that claims a universal principle and validates it across multiple
domains must be explicit about the boundary between proven and projected.
This chapter draws that boundary for the ORIUS framework.

\subsection{The Full Formal Apparatus}

The formal surface of this dissertation is larger than the T1--T8 label
suggests.  The integrated audit (Table~\ref{tab:integrated-theorem-surface-summary})
records \textbf{82 formal items}: 26~theorems, 5~lemmas, 8~propositions,
12~corollaries, 15~definitions, and 16~standing assumptions.  The active
theorem count is computed over the master-manuscript include tree rather
than the entire repository, and all 26 theorems pass the integrated theorem
gate with locked code anchors and
artifact traceability.

The foundational theorems precede the formal DC3S ladder:

\begin{description}
  \item[\textbf{S1} (Existence of the Illusion under Dropout).]
    Established in Chapter~\ref{ch:safety-illusion},
    S1 proves that under any positive dropout rate a deterministic
    controller can satisfy its observed constraint while violating the
    true physical envelope.  This is the phenomenon that motivates
    the entire framework.
  \item[\textbf{S2} (Informal Safety Guarantee of DC3S).]
    Also in Chapter~\ref{ch:safety-illusion}, S2 gives
    the informal correctness claim: DC3S closes the gap S1 identifies
    by inflating the feasible action set by the reliability-adjusted
    conformal quantile.
\end{description}

The eight principal theorems T1--T8 then formalise, quantify, and
extend this guarantee across one-step, episode-level, and multi-step
temporal horizons.

\subsection{What Is Fully Validated Across the Tiered Runtime Package}

The following results are established by reproducible experiment on the locked
CPSBench harness (seed~42, $N = 2000$ fault-stressed steps, unified fault schedule:
dropout~15\%, spike~8\%, stale~10\%).  The formal reference surface remains
battery-first, and the cross-domain claim strength is governed by the
reference / proof-validated / proof-candidate / shadow-synthetic /
experimental tier matrix:

\begin{enumerate}
  \item \textbf{OASG existence.}  The Observation-Action Safety Gap is
    measurably non-zero across the locked runtime package under the shared
    fault schedule.  Baseline TSVR is non-zero in the reference row, both
    proof-validated peers, the AV proof-candidate row, the aerospace
    experimental row, and the navigation shadow-synthetic row.
  \item \textbf{DC3S effectiveness.}  The DC3S shield achieves
    0.0\% TSVR across the three proof-validated domains (energy management,
    industrial process control, and medical monitoring) under the locked
    three-seed, 48-step replay protocol.  AV and aerospace show TSVR
    reduction under the higher-budget N~$= 2000$ evaluation and are
    classified as proof-candidate and experimental surfaces respectively.
  \item \textbf{Zero true-state violation rate.}  Energy management,
    industrial, and healthcare each achieve 0.0\% TSVR under DC3S in the
    proof evidence package.  AV remains proof-candidate and aerospace
    remains experimental even where higher-budget replay shows additional
    repair benefit.
  \item \textbf{Formal safety bound.}  Theorem~\ref{thm:safety-preservation}
    guarantees that the true-state violation probability is at most $\alpha$
    with confidence $1 - \delta$, under exchangeability of the calibration
    residuals and bounded observation error.  The bound is formalized on the
    battery reference surface and instantiated across the tiered replay rows
    via the shared domain-adapter mapping.
  \item \textbf{Bounded intervention cost.}  The shield's repair rate remains
    below 100\% across the five forecasting-capable runtime rows (mean~$\bar{w}_t \in
    [0.66, 0.81]$), confirming that DC3S does not collapse all actions to a
    constant fallback.
  \item \textbf{Sub-millisecond overhead.}  Pipeline latency (P95~$<$~0.1~ms)
    is consistent across the five forecasting-capable runtime rows, confirming runtime
    suitability for fast control-loop applications.
\end{enumerate}

Navigation remains simulation-backed portability evidence outside this
forecasting-surface summary.

\subsection{What Is Fully Validated in Battery (Reference Domain)}

The battery domain provides the complete theorem-to-code-to-artifact lineage:

\begin{enumerate}
  \item \textbf{OASG existence.}  Measurably non-zero at dropout rates
    as low as~5\%.  Chapter~\ref{ch:main-results} reports TSVR of
    3.9--4.2\% for the deterministic LP baseline at 20\% dropout.
  \item \textbf{DC3S effectiveness.}  Exactly 0\% TSVR across all tested
    dropout rates $d \in [0, 0.30]$, both grid regions, and all fault combinations.
  \item \textbf{Full formal proof surface.}  Theorems T1--T8 are stated,
    proved, and connected to code artifacts for the battery domain.
    Chapter~16--20 each carry domain-instantiation subsections extending
    each theorem to the peer CPS domains.
  \item \textbf{Intervention cost bounded.}  Repair rate remains below
    3.1\% under \texttt{drift\_combo} and below 20\% even at 30\% dropout.
\end{enumerate}

\subsection{What Is Partially Proved}

Several claims are supported by evidence but carry caveats:

\begin{enumerate}
  \item \textbf{Conditional coverage.}  CQR group coverage
    (Chapter~\ref{ch:conditional-coverage}) is validated on the four
    tested fault families.  Coverage under \emph{out-of-distribution}
    faults not in the calibration set has not been verified.
  \item \textbf{Aging effects.}  Aging-aware recalibration
    (Chapter~\ref{ch:temporal-extensions}) is demonstrated on synthetic
    degradation curves.  Multi-year field aging data is not available.
  \item \textbf{Certificate half-life.}  The blackout study
    (Chapter~\ref{ch:cert-half-life-blackout}) establishes a bounded
    safe-hold result; longer horizons (weeks, months) remain extrapolated.
  \item \textbf{Transfer across regions.}  The US--DE transfer stress
    test (Chapter~\ref{ch:transfer-generalisation}) shows positive
    results for two grid regions.  Grids with fundamentally different
    market structures are untested.
\end{enumerate}

\subsection{Why Multi-Domain Validation Strengthens the Core Claim}

The universal validation is not scope-widening---it is confirmation that
the core claim (DC3S closes the OASG) is not an artefact of battery physics.
Each additional domain that passes the same promotion gate (locked telemetry,
verified training surface, SIL trace pass, material TSVR reduction) provides
independent evidence that the OASG is a structural phenomenon in any CPS with
degraded telemetry, not a battery-specific measurement anomaly.

Energy management provides the canonical instantiation with full theorem lineage.
Industrial and healthcare provide proof-validated empirical depth under the
shared universal replay harness.  AV contributes proof-candidate evidence at
higher evaluation budgets, aerospace contributes experimental evidence, and
navigation provides simulation-backed
portability evidence.  Together, the five real-data runtime rows provide
broad evidence of the OASG phenomenon and DC3S effectiveness under the ORIUS
principle, while the final claim boundary still follows the locked tier
matrix rather than a flat reading of every row.

\subsection{What Remains Outside the Current Evidence}

\begin{enumerate}
  \item \textbf{Field deployment.}  No domain has been validated on production
    hardware in a live safety-critical loop.
  \item \textbf{Heterogeneous multi-asset composition.}  Fleet-level shielding
    is validated for a two-battery shared-transformer case; cross-domain
    composition (battery + AV + medical device) is an open problem.
  \item \textbf{Adversarial completeness.}  The active-probing audit detects
    the tested spoofing strategy.  Completeness against all adversarial
    strategies is not claimed.
  \item \textbf{Non-box constraint geometry.}  All current domain adapters
    use box or interval constraints.  Polytopic and conic constraint
    structures are architectural claims, not yet empirically validated.
\end{enumerate}

\subsection{Chapter Summary}

This thesis contributes a 67-item formal apparatus (18 theorems, 5 lemmas,
7 propositions, 10 corollaries, 11 definitions, 16 assumptions) grounded
in 25 chapter-level traceability rows connecting each formal result to its
code anchor and artifact evidence.  The ORIUS principle is established on a
tiered runtime package: the OASG exists across the locked runtime rows,
DC3S achieves zero TSVR in the three highest-confidence rows (battery
reference, industrial, healthcare), AV remains proof-candidate, aerospace
remains experimental, and navigation adds simulation-backed portability
evidence.  Being explicit about what remains outside the battery reference
row and proof-validated peer rows
protects the core multi-domain result from credibility it does not yet claim.


% --- Merged from ch33_what_battery_thesis_proves.tex ---
\section{What the ORIUS Framework Proves}
\label{ch:what-battery-thesis-proves}

This chapter draws a strict line between what the ORIUS evidence package
establishes and what remains outside its current scope.  The value of the
framework is not breadth of claim, but depth of traceability: every asserted
result connects from theorem statement through code to artifact, with the
strength of each cross-domain claim determined by its evidence tier.

\subsection{What Is Established}

The ORIUS framework proves the following, across the defended evidence package:

\begin{enumerate}
  \item \textbf{OASG existence is universal.}  The observation-action safety
    gap is not domain-specific.  It arises whenever a closed-loop controller
    optimises against a degraded observation channel.  Evidence of the gap
    was measured across energy management, autonomous vehicles, industrial
    process control, medical monitoring, aerospace, and navigation.

  \item \textbf{DC3S closes the gap reliably.}  In the three
    proof-validated domains (energy management, industrial process control,
    and medical monitoring), DC3S reduces the true-state violation rate
    to zero (TSVR~$= 0$\%) while maintaining dispatch performance comparable
    to the unshielded controller.  AV carries proof-candidate evidence and
    aerospace carries experimental evidence of material TSVR reduction under
    the higher-budget N~$= 2000$ evaluation
    but are classified as proof-candidate and experimental surfaces respectively
    under the locked three-seed, 48-step protocol.

  \item \textbf{The safety bound is quantitative and verifiable.}
    Theorem~\ref{thm:core-bound} establishes
    $\mathbb{E}[V] \le \alpha(1-\bar{w})T$ under Assumptions A1--A7.  The
    finite-sample correction $\varepsilon_n = O(1/\sqrt{n})$ tightens the
    bound for calibration sets of finite size.  Every empirical TSVR lies
    strictly below this bound with measurable slack.

  \item \textbf{Traceability is end-to-end.}  Each theorem group (T1--T8)
    is linked to a code anchor in \texttt{src/orius/dc3s/} and to a locked
    artifact trace.  The integrated theorem surface register
    (Appendix~\ref{app:integrated-theorem-surface-register}) maps claim to code to artifact
    for every defended result.

  \item \textbf{The hardening surface is deep.}  Beyond the core TSVR
    result, the energy management domain contributes: fault-performance
    stress tests, latency benchmarking (P95 $< 5$~ms end-to-end),
    conditional coverage audits, hyperparameter stability analysis,
    hardware-in-the-loop rehearsal, certificate half-life studies, graceful
    degradation protocols, two-asset fleet compositional safety, and
    adversarial robustness via active probing.
\end{enumerate}

\subsection{Theorem-to-Empirical Correspondence}

Table~\ref{tab:theorem-empirical-correspondence} maps the four core theorems
to their empirical confirmations across domains.

\begin{table}[htbp]
\centering
\caption{Core theorem ladder and empirical confirmation across the tiered six-row runtime package.}
\label{tab:theorem-empirical-correspondence}
\begin{tabular}{p{3.8cm}p{5.5cm}p{4.5cm}}
\toprule
\textbf{Theorem} & \textbf{Prediction} & \textbf{Empirical confirmation} \\
\midrule
T1 -- OASG Existence
  & TSVR $> 0$ under quality-ignorant controller
  & 2.78\%--21.53\% baseline TSVR across the locked non-battery runtime rows \\
\addlinespace
T2 -- Safety Preservation
  & $\Pr[\phi_{\text{true}}(a_t) = 1] \ge 1-\alpha$ per step
  & TSVR $= 0$\% across the battery reference row and the two proof-validated peers \\
\addlinespace
T3 -- Core Bound
  & $\mathbb{E}[V] \le \alpha(1-\bar{w})T$
  & Runtime TSVRs stay within the conservative battery-anchored envelope \\
\addlinespace
T4 -- No-Free-Safety
  & Any safe-in-observation-only policy has TSVR $> 0$
  & Rule-based baselines confirm non-zero TSVR on the locked runtime rows \\
\bottomrule
\end{tabular}
\end{table}

\subsection{What the Framework Result Means}

The ORIUS result is strong because it ties every major claim to an operational
consequence.  Hidden violations are not abstract theorem objects; they
correspond to unsafe physical outcomes---damaging battery discharge, vehicle
over-speed, process power exceedance, ICU alarm suppression failure, or
aircraft envelope breach.  Reliability-aware interval inflation is not only a
calibration device; it widens the safe-action margin in proportion to the
observation degradation severity.  Action repair is not merely a projection
lemma; it is the mechanism that prevents unsafe actuation when the observed
state is too optimistic.  Certificates are not audit theatre; they provide
a per-step commitment to the safe-action guarantee that can be logged,
reviewed, and acted upon by downstream safety monitors.

\subsection{Scope of the Final Claim}

The final claim is precisely bounded:

\begin{quote}
Under degraded telemetry, a quality-ignorant controller can be observationally
safe while physically unsafe.  The DC3S shield addresses this gap through
reliability-aware uncertainty inflation, safe-action repair, and auditable
runtime certificates.  In this thesis, the complete formal reference
lineage is battery; industrial and healthcare are proof-validated empirical
peers; AV remains proof-candidate; aerospace remains experimental; and
navigation remains shadow-synthetic portability evidence.
\end{quote}

This claim is supported by code, artifacts, and chapter-level alignment.  It
does not extend to production deployment, adversarial completeness, or
heterogeneous multi-domain composition without further engineering validation
(Chapter~\ref{ch:outside-current-evidence}).

\subsection{Summary}

The ORIUS framework proves that the observation-action safety gap is a
universal CPS phenomenon with a principled, deployable remedy.  The energy
management domain provides the formally locked reference surface; four
additional domains confirm the result with equivalent replay depth.  The
theorem ladder, finite-sample bound, and cross-domain PICP confirmation
together constitute a complete, end-to-end safety argument for DC3S-shielded
controllers under degraded telemetry.


% --- Merged from ch35_deployment_path_verification_discipline.tex ---
\section{Deployment Path and Verification Discipline}
\label{ch:deployment-path-verification}

The ORIUS framework closes with a deployment path that is operationally
serious but evidence-bounded.  The goal is not to announce production
deployment.  The goal is to specify how deployment would proceed for any
domain without breaking the claim discipline established in this thesis.

\subsection{Universal Deployment Readiness Checklist}

Before deploying DC3S in any domain, the following conditions must be
verified:

\begin{enumerate}
  \item \textbf{Domain adapter registration.}  A concrete \texttt{DomainAdapter}
    subclass implementing all five pipeline stages (ingest, OQE, uncertainty,
    tighten, repair, certify) must be registered in the domain registry.
    All adapter unit tests must pass.

  \item \textbf{Calibration surface validation.}  A domain-specific calibration
    split must achieve $\mathrm{PICP}_{90} \ge 0.90$.  The calibration
    residuals must satisfy approximate exchangeability with the deployment
    telemetry distribution.

  \item \textbf{OQE threshold configuration.}  The reliability score $w_t$
    thresholds (alert level, blackout level) must be calibrated against
    domain-specific telemetry fault rates.  Stale, spike, and dropout
    fault injection tests must confirm OQE detection lag within the
    Assumption~A6 bound.

  \item \textbf{Constraint and feasibility verification.}  The safety predicate
    $\phi(x_t)$, the constraint set $\mathcal{U}$, and the margin function
    $m_t = q_t / (w_t + \varepsilon)$ must be reviewed by domain experts.
    Assumption~A3 (feasible safe repair) must be confirmed: $\mathcal{A}_t$
    must be non-empty whenever the domain is in a safe physical state.

  \item \textbf{End-to-end benchmark replay.}  A CPSBench replay with separate
    truth and observed trajectories must be executed over at least 1000
    telemetry steps under the target fault schedule.  The empirical TSVR
    must satisfy the Core Bound ($\mathbb{E}[V] \le \alpha(1-\bar{w})T$).

  \item \textbf{Artifact manifest traceability.}  Every output artifact
    (tables, figures, certificates) must be registered in the artifact
    manifest with its generating script, run parameters, and version hash.
    Re-running the named script must reproduce the artifact within numerical
    tolerance.
\end{enumerate}

\subsection{Universal Deployment Ladder}

For any CPS domain, the deployment ladder is:

\begin{enumerate}
  \item \textbf{Benchmark replay}---observed versus true trajectory separation,
    TSVR measurement, Core Bound verification.
  \item \textbf{Software-in-loop (SIL) rehearsal}---DC3S running against a
    simulated or replayed telemetry stream, with full certificate emission
    and manifest logging.
  \item \textbf{Hardware-in-loop (HIL) validation}---real telemetry transport,
    physical actuator interface, relay behavior under injected faults.
  \item \textbf{Supervised pilot operation}---bounded fallback policies,
    certificate-gated alert suppression, human-in-the-loop audit.
  \item \textbf{Wider deployment}---only after all prior layers remain
    manifest-traceable and no regression in TSVR or certificate integrity
    has been observed.
\end{enumerate}

The energy management domain is currently at Stage~1--3.  No domain in
this thesis has reached Stage~4.  Stages~4--5 require domain-specific
regulatory engagement beyond the scope of this framework contribution.

\subsection{Operational Monitoring}

Once deployed in any live tier, DC3S requires continuous monitoring of:

\begin{itemize}
  \item \textbf{OQE trend.}  A sustained drop in the average reliability
    score $\bar{w}_t$ below the alert threshold ($w < 0.60$ for energy
    management, domain-specific for others) should trigger a calibration
    review and possible shield recalibration.

  \item \textbf{Intervention rate drift.}  An intervention rate (IR) that
    exceeds the baseline by more than 10 percentage points may signal
    distribution shift in the input telemetry.  The conformal quantile
    $q_t$ should be re-estimated on a fresh calibration split.

  \item \textbf{Certificate half-life.}  The runtime certificate validity
    window should be monitored using the decay model in
    Chapter~\ref{ch:cert-half-life-blackout}.  Certificates expiring faster
    than expected indicate elevated fault rates or OQE degradation.

  \item \textbf{TSVR monitoring.}  Where ground-truth state can be accessed
    post-hoc (e.g., from SCADA historian data or sensor reconciliation),
    the empirical TSVR should be computed at each reporting period and
    compared to the Core Bound prediction.
\end{itemize}

\subsection{Rollback Protocol}

If any monitoring indicator exceeds its alarm threshold, the following
rollback protocol applies:

\begin{enumerate}
  \item \textbf{Certificate hold.}  Suspend certificate emission and fall
    back to the safe-landing policy (Chapter~\ref{ch:graceful-degradation}).
  \item \textbf{Fault diagnosis.}  Identify the root cause of the degradation
    (telemetry faults, model drift, constraint boundary change) using the
    OQE diagnostic flags.
  \item \textbf{Calibration refresh.}  Re-run conformal calibration on a
    fresh holdout split reflecting current operating conditions.
  \item \textbf{Re-verification.}  Execute a full SIL pass under the refreshed
    calibration before resuming certificate emission.
\end{enumerate}

\subsection{Verification Discipline}

Every stage of the deployment ladder must preserve four invariants:

\begin{itemize}
  \item a named script path that can regenerate each output artifact,
  \item a run-scoped artifact with stable, version-controlled location,
  \item a manifest entry tying the artifact to the corresponding claim, and
  \item an explicit statement of what the result does \emph{not} prove.
\end{itemize}

This discipline is not only an engineering convenience.  It is a
prerequisite for any safety-critical certification process.  A domain
that cannot reproduce its validation artifacts on demand cannot make a
defensible safety argument to a regulatory body.

\subsection{Summary}

The ORIUS deployment path is incremental and auditable.  Each new layer
inherits the theorem-to-code-to-artifact traceability that governs the
current evidence package.  Widening to production---within any domain---is
a matter of climbing the deployment ladder step by step, with each rung
verified against the Core Bound and logged in the artifact manifest.


% --- Merged from ch23_research_roadmap.tex ---

\section{Research Roadmap and Deployment Path}
\label{ch:roadmap}

This chapter lays out the sequence of work required to move ORIUS from a
battery-validated prototype to a multi-domain, deployable safety runtime.
Each milestone is ordered by dependency: later items require earlier ones,
and no step can be skipped without accumulating technical debt.

\subsection{Submit the Battery Foundation First}

The immediate priority is to publish the battery-domain results as a
self-contained contribution.  The submission package comprises the CPSBench
harness, the DC3S implementation, the dual-trajectory evaluation protocol,
and the empirical evidence of zero true-state violations.  Peer review of
the battery result establishes the OASG concept and the DC3S shield in the
literature before generalisation claims are made.

\subsection{Fix Remaining Empirical Gaps}

Before extending to new domains, the battery evidence must be hardened:

\begin{enumerate}
  \item \textbf{Out-of-distribution faults.}  Extend the fault model to
    include correlated multi-sensor dropout and Byzantine sensor failures.
    Validate that CQR residuals remain exchangeable under these conditions
    or document the breakdown.
  \item \textbf{Long-horizon aging.}  Obtain or simulate multi-year
    degradation trajectories and verify that the aging-aware recalibration
    of Chapter~\ref{ch:temporal-extensions} maintains coverage.
  \item \textbf{Regional breadth.}  Run the full battery on at least two
    additional grid regions (e.g., ERCOT, NordPool) to confirm that
    market-structure differences do not invalidate the safety guarantee.
\end{enumerate}

\subsection{Extract ORIUS Core Modules}

Refactor the current codebase to separate domain-independent modules
(OQE scoring, RAC inflation, certificate issuance, audit trail) from
battery-specific code (SOC dynamics, battery constraints, SCADA parser).
The result is a \texttt{orius-core} library with a stable API and a
\texttt{orius-battery} adapter that depends on it.  The separation must
be enforced by interface boundaries, not merely by convention.

\subsection{Build the Battery Adapter Package}

Package the battery-specific adapter as a standalone installable module
that depends on \texttt{orius-core}.  This package includes:

\begin{itemize}
  \item The SOC integrator plant model.
  \item The SCADA telemetry parser and OQE feature extractor.
  \item The box-constraint projection oracle.
  \item The CPSBench battery plant for regression testing.
\end{itemize}

Publishing this package demonstrates the adapter pattern in practice and
provides a template for subsequent domain adapters.

\subsection{Add a Second Domain}

The second domain should stress the aspects of ORIUS that battery did not
test: nonlinear dynamics, multi-dimensional state, and non-box constraints.
Candidates include building HVAC (nonlinear thermal dynamics, polytopic
comfort constraints) and autonomous vehicle lane-keeping (six-DOF state,
curvature-dependent road constraints).  The second-domain study must
reproduce the full evaluation protocol: dual-trajectory simulation, TSVR
measurement, conformal coverage audit, and intervention-rate analysis.

\subsection{Add a Third Domain and Beyond}

Each additional domain further validates the invariant core and reveals
edge cases in the adapter interface.  After the third domain the
\texttt{orius-core} API should be considered stable; subsequent domains
become engineering exercises rather than research contributions.

\subsection{Hardware-in-the-Loop Validation}

The CPSBench harness simulates telemetry faults in software.  Before
deployment, the DC3S shield must be tested with \emph{real} packet loss,
delay jitter, and sensor degradation on a hardware-in-the-loop (HIL)
testbed.  The HIL stage validates:

\begin{enumerate}
  \item That the OQE feature extractor behaves correctly on raw network
    packets, not simulated dropout masks.
  \item That the shield's latency budget (established in
    Chapter~\ref{ch:latency}) holds on embedded hardware.
  \item That the certificate chain persists correctly under power
    interruptions and storage faults.
\end{enumerate}

\subsection{ORIUS-Bench Public Release}

The CPSBench harness should be released as a public benchmark---ORIUS-Bench
---with a standardised evaluation API, a leaderboard, and reference
implementations for at least two domains.  Public benchmarking invites
independent reproduction and adversarial stress-testing by the community,
which no internal audit can substitute.

\subsection{CertOS Runtime Path}

The long-term deployment target is a \emph{CertOS}: a lightweight runtime
operating system that wraps any black-box controller with the ORIUS
safety layer.  CertOS would run as a sidecar process, intercept actuation
commands, evaluate them against the current certificate, and either pass
or project.  The runtime must satisfy hard real-time constraints, which
imposes implementation requirements beyond the Python prototype used in
this dissertation.

\subsection{Chapter Summary}

The roadmap proceeds in strict dependency order: publish the battery
foundation, harden its empirical gaps, extract the domain-independent
core, validate on a second and third domain, run hardware-in-the-loop
tests, release ORIUS-Bench publicly, and ultimately build a deployable
CertOS runtime.  Each step is necessary; none is sufficient alone.  The
battery result is the foundation, not the destination.


% --- Merged from ch36_conclusion.tex ---
\section{Conclusion}
\label{ch:conclusion}

The central finding of this thesis is concrete: when telemetry degrades,
a controller can satisfy its observed constraints while violating the true
physical envelope of the system.  This hidden gap---the
Observation-Action Safety Gap (OASG)---is measurable, operationally
meaningful, and correctable.  ORIUS demonstrates this across a tiered
six-row runtime package.  The battery reference row and the two
proof-validated peers (industrial process control and medical monitoring)
achieve zero true-state violations under DC3S in the locked three-seed,
48-step protocol; autonomous vehicles remains proof-candidate, aerospace
remains experimental, and navigation provides simulation-backed portability
evidence.

\subsection{Universal Framework Outcome}

The ORIUS framework demonstrates that the OASG phenomenon is not a
battery-specific artifact.  It is a structural consequence of any
closed-loop controller that optimises against degraded telemetry.  The
five-stage DC3S pipeline---Detect, Calibrate, Constrain, Shield,
Certify---operates on domain-abstract objects (observation vector,
reliability score, conformal region, feasible set, projected action) and
therefore transfers to any CPS for which a \texttt{DomainAdapter} can be
implemented.

\subsection{Domain Coverage and Empirical Results}

The six-row runtime package was exercised under a unified fault schedule
(dropout~15\%, spike~8\%, stale~10\%) at $N = 2000$ fault-stressed steps:

\begin{itemize}
  \item \textbf{Energy Management} (reference): 3.9\% baseline TSVR $\rightarrow$
    0.0\% under DC3S (locked; proof-validated).  Full theorem-to-code-to-artifact
    lineage established in Parts~I--II.
  \item \textbf{Industrial Process Control}: 21.53\% $\rightarrow$ 0.0\%
    (91.7\% repair rate, locked; proof-validated).
  \item \textbf{Medical Monitoring (ICU)}: 6.25\% $\rightarrow$ 0.0\%
    (14.6\% repair rate, locked; proof-validated).
  \item \textbf{Autonomous Vehicles}: 2.78\% $\rightarrow$ 2.78\%
    (98.6\% repair rate, locked; proof-candidate---TSVR unchanged under
    locked 48-step protocol due to constraint structure; material reduction
    observed at N~$= 2000$ evaluation budget).
  \item \textbf{Aerospace (Flight Envelope)}: 9.72\% $\rightarrow$ 9.72\%
    (13.9\% repair rate, locked; experimental surface---TSVR unchanged
    under locked 48-step protocol; further evaluation ongoing).
\end{itemize}

Navigation remains part of the broader ORIUS portability package, but it is
simulation-backed rather than part of the proof-validated claim set:
25.8\% $\rightarrow$ 9.25\% (100\% repair rate, mean $\bar{w}_t = 0.773$).
Pipeline latency remains sub-millisecond across the tiered runtime package
(P95~$<$~0.1~ms), confirming ORIUS is suitable for fast-control-loop
applications on the replay surfaces evaluated here.

\subsection{Formal Apparatus and Theorem Ladder Scope}

The full formal apparatus of this dissertation consists of 26 theorems,
5 lemmas, 8 propositions, 12 corollaries, 15 definitions, and 16 standing
assumptions---82 formal items in total.  Two foundational theorems are
established in Part~I before the formal ladder:

\begin{itemize}
  \item \textbf{S1} (Theorem: Existence of the Illusion under Dropout,
    Chapter~\ref{ch:safety-illusion}): proves that a
    controller can satisfy observed constraints while violating true
    physical safety under any positive dropout rate.
  \item \textbf{S2} (Theorem: Informal Safety Guarantee of DC3S,
    Chapter~\ref{ch:safety-illusion}): establishes the
    informal correctness guarantee that motivates the DC3S design.
\end{itemize}

The eight principal theorems of the formal DC3S ladder (T1--T8) are then
stated, proved, and connected to code artifacts in
Chapters~\ref{ch:theorem-oasg}--\ref{ch:temporal-extensions}.
The Domain Instantiation subsections in each of those chapters demonstrate
that no battery-specific assumption is required: each abstract object
($z_t$, $w_t$, $C_t^{\text{RAC}}$, $\mathcal{A}_t$, $a_t^{\text{safe}}$,
$v(\cdot)$) maps directly to its domain-specific counterpart in AV,
industrial, healthcare, and aerospace.  All 18 theorems pass the integrated
theorem gate with locked code anchors and artifact traceability
(see \texttt{reports/publication/integrated\_theorem\_gate.csv}).

\subsection{Final Boundary}

This thesis does not claim that ORIUS is ready for deployment in any
safety-critical domain without domain-specific validation.  It claims
that:

\begin{enumerate}
  \item The OASG phenomenon is real and measurable in any CPS with
    degraded telemetry.
  \item The DC3S architecture provides a principled correction mechanism
    that is domain-agnostic by construction.
  \item Shared replay evaluation shows consistent runtime behavior across
    the six-row package, while claim strength follows the locked tier table
    rather than raw replay metrics alone.
  \item The energy management domain provides the complete reference
    instantiation with full theorem-to-artifact lineage; industrial and
    healthcare provide proof-validated empirical confirmation; AV is
    proof-candidate; aerospace carries experimental evidence at additional
    evaluation budgets;
    and navigation provides simulation-backed portability evidence.
\end{enumerate}

\subsection{Closing Statement}

The ORIUS thesis ends where a rigorous framework contribution should: with
a locked, traceable, reproducible tiered multi-domain safety package, with
battery as the formal reference row, industrial and healthcare as the
proof-validated peers, AV as proof-candidate, aerospace as experimental,
navigation as simulation-backed portability evidence, and deeper
domain-specific validation still required at each deployment site.


% --- Merged from ch24_conclusion.tex ---

\section{Synthesis and Outlook}
\label{ch:synthesis-outlook}

Physical AI controllers routinely make decisions based on observations they
have no reason to trust.  A telemetry packet arrives late, or not at all,
or carrying a stale reading---and the controller evaluates its safety
constraints against a state the plant may never have occupied.  The
resulting Observation-Action Safety Gap is invisible to any evaluation
protocol that checks constraints only on the observed trajectory.

This dissertation has shown, in the battery dispatch domain, that the gap
is real, measurable, and dangerous.  More importantly, it has shown that
the gap is \emph{closable}.  The ORIUS principle---make observation
reliability part of the control state, inflate uncertainty in proportion to
unreliability, and project actions onto the tightened feasible set---eliminates
true-state violations without replacing the upstream optimiser or requiring
a new plant model.

Under degraded observation, safety cannot be taken for granted.  ORIUS makes
observation reliability a first-class quantity in the control loop, not an
afterthought logged for post-hoc analysis.

\subsection{What the Battery Domain Taught Us}

Grid-scale battery dispatch was chosen as the first proving ground because
it combines hard state constraints, high action-to-state sensitivity, and
network-mediated telemetry.  The experimental campaign yielded several
lessons that extend beyond batteries:

\begin{enumerate}
  \item \textbf{Hidden violations are common.}  Under 20\% dropout, a
    deterministic LP violates true-state SOC bounds at 3.9--4.2\% of
    steps while reporting full feasibility.  The controller believes it
    is safe; the plant disagrees.
  \item \textbf{Conformal prediction is the right uncertainty language.}
    Distribution-free coverage guarantees allow the safety layer to operate
    without a parametric model of telemetry faults.  The CQR intervals
    adapt to the empirical residual distribution, including heavy tails
    induced by sensor drift.
  \item \textbf{Reliability-adaptive inflation is the key mechanism.}
    A conformal interval of fixed width is either too wide when telemetry
    is good (excessive conservatism) or too narrow when telemetry is
    degraded (missed violations).  Scaling width by $1/w_t$ tracks the
    actual risk, concentrating intervention in the episodes that need it.
  \item \textbf{The cost of safety is modest.}  Intervention rates of
    2.8--3.1\% under the canonical fault scenario translate to fewer than
    three corrections per hundred dispatch intervals---a negligible
    operational burden for the elimination of all hidden violations.
\end{enumerate}

\subsection{Why Observation Reliability Belongs Inside Control}

Classical control treats observation noise as an exogenous disturbance to
be filtered, not a quantity the controller should reason about.  This
separation is justified when the noise model is accurate and stationary.
In modern CPS, neither condition holds: telemetry faults are non-stationary,
packet-level, and correlated with system load.  Treating reliability as
external means the controller cannot tighten its actions when observation
quality degrades, and it cannot relax them when quality is high.

ORIUS brings reliability inside the loop.  The OQE score $w_t$ is updated
at every step from raw telemetry features.  It feeds directly into the
uncertainty inflator, which sets the width of the safe-action region.  The
controller's conservatism is therefore \emph{proportional to its ignorance}:
when it knows little, it acts cautiously; when it knows much, it acts
freely.  This proportionality is the architectural insight of the
dissertation, and it is domain-independent.

\subsection{What ORIUS Can Become}

The battery result proves the concept.  The adapter architecture
(Chapter~\ref{ch:battery-to-orius}) provides the blueprint for
generalisation.  If the ORIUS community delivers on the roadmap of
Chapter~\ref{ch:roadmap}---a second domain, a public benchmark,
a deployable runtime---the result is a certification layer that any
cyber-physical controller can use, regardless of its internal optimisation
logic.

This is not a small ambition.  It requires that the conformal coverage
guarantee survives domain transfer, that the adapter interface is
expressive enough for nonlinear plants with high-dimensional state, and
that the runtime meets hard real-time constraints on embedded hardware.
Each of these is a non-trivial research or engineering challenge.  But
the structural foundations---domain-abstract scoring, distribution-free
uncertainty, worst-case projection---are in place.

\subsection{Final Remarks}

This dissertation began with a question: what happens when a controller
cannot trust its own eyes?  The answer, demonstrated rigorously in one
domain and argued structurally for others, is that safety must be
\emph{earned} at every time step by measuring observation quality,
quantifying the resulting uncertainty, and restricting actions to those
that remain safe under the worst case.

The battery domain was the right place to start.  It will not be the
last place ORIUS is needed.
